% ============================================================
% CAPÍTULO 4: RESULTADOS
% ============================================================

\chapter{Resultados y Discusión}
\label{chap:resultados}

En este capítulo se presentan los resultados obtenidos por las cuatro
arquitecturas implementadas (ResNet 1D, SEResNet 1D, el modelo
híbrido CNN-Transformer y Mamba 1D). El entrenamiento se ha realizado mediante
una validación cruzada por grupos (5-Fold Group Cross-Validation),
asegurando que los registros de un mismo sujeto
nunca se compartan entre los conjuntos de entrenamiento y validación. 

Para la comparación con otros trabajos publicados, también se evalua en un conjunto
de test independiente ~\cite{moody_paf_2001}.

\section{Resultados detallados por arquitectura}
\label{sec:resultados_por_arquitectura}

A continuación se presentan los resultados obtenidos para cada una de las arquitecturas.

\subsection{Arquitectura ResNet 1D}
\label{subsec:resultados_resnet}

Para optimizar el rendimiento de la ResNet 1D,
se parametrizó la arquitectura base para y evaluar el impacto
de las siguientes variables de diseño:
\begin{itemize}
    \item \textbf{Profundidad de la red}: Variación del número de bloques residuales por capa.
    El baseline original utiliza 8 bloques (configuración \texttt{[2, 2, 2, 2]}), mientras que la variante
    profunda utiliza 12 bloques (\texttt{[3, 3, 3, 3]}).
    \item \textbf{Ancho de la red}: Modificación del número de canales convolucionales en las capas sucesivas ( baseline original: \texttt{[32, 64, 128, 256]} frente a la versión ancha: \texttt{[64, 128, 256, 512]}).
    \item \textbf{Tamaño del kernel}: Comparación del tamaño del kernel en las convoluciones 1D (baseline: 15 muestras, frente a kernels grandes de 23 y pequeños de 7).
    \item \textbf{Profundidad y Estrechez}: Una combinación de mayor profundidad (15 bloques distribuidos como \texttt{[3, 4, 6, 3]}) con un menor número de canales (\texttt{[24, 48, 96, 192]}) para restringir el número total de parámetros de la red y mitigar el sobreajuste.
    \item \textbf{Aumento de Datos (Data Augmentation)}: Evaluación de cada una de las arquitecturas anteriores bajo aumentación de datos (ruido gaussiano, baseline wander y escalado de señal).
\end{itemize}

Los resultados obtenidos por cada variante mediante validación cruzada de
5 \textit{folds} por grupos se resumen en la
Tabla~\ref{tab:resultados_resnet_variaciones}.

\begin{table}[htbp]
  \centering
  \caption{Resultados de las variaciones de la arquitectura ResNet 1D con y sin aumento de datos}
  \label{tab:resultados_resnet_variaciones}
  \resizebox{\textwidth}{!}{%
  \begin{tabular}{llcccc}
    \toprule
    \textbf{Configuración ResNet 1D} & \textbf{Augment.} & \textbf{F1 Macro Val (CV)} & \textbf{F1 Macro Test} & \textbf{Accuracy Test} & \textbf{Score Reto (Event 2)} \\
    \midrule
    \multirow{2}{*}{ResNet Base} & No & $0.7591 \pm 0.04$ & 0.80 & 0.80 & 46\% \\
                                 & Sí & $0.7473 \pm 0.04$ & 0.78 & 0.79 & \textbf{64\%} \\
    \midrule
    \multirow{2}{*}{\textbf{ResNet Profunda (\textit{deep})}} & No & $\mathbf{0.7863 \pm 0.04}$ & 0.82 & 0.82 & 36\% \\
                                                     & Sí & $0.7501 \pm 0.04$ & 0.80 & 0.80 & 50\% \\
    \midrule
    \multirow{2}{*}{ResNet Ancha (\textit{wide})} & No & $0.7439 \pm 0.04$ & 0.77 & 0.77 & 50\% \\
                                                  & Sí & $0.7447 \pm 0.06$ & 0.80 & 0.81 & 57\% \\
    \midrule
    \multirow{2}{*}{ResNet Kernel Grande (k=23)} & No & $0.7464 \pm 0.04$ & 0.80 & 0.80 & 54\% \\
                                                  & Sí & $0.7417 \pm 0.05$ & 0.77 & 0.77 & 46\% \\
    \midrule
    \multirow{2}{*}{ResNet Kernel Pequeño (k=7)} & No & $0.7378 \pm 0.06$ & \textbf{0.83} & \textbf{0.83} & 57\% \\
                                                  & Sí & $0.7461 \pm 0.04$ & 0.79 & 0.79 & 46\% \\
    \midrule
    \multirow{2}{*}{ResNet Profunda y Estrecha} & No & $0.7496 \pm 0.06$ & 0.81 & 0.81 & 54\% \\
                                                         & Sí & $0.7306 \pm 0.06$ & 0.79 & 0.80 & 54\% \\
    \bottomrule
  \end{tabular}%
  }
  \vskip 5pt
  \small
  \textit{Nota:} El score del reto representa la evaluación en el conjunto de test blindado de afpdb~\cite{moody_paf_2001}.
\end{table}

De los resultados de la Tabla~\ref{tab:resultados_resnet_variaciones}, vemos como en este caso La ResNet Profunda (deep) sin aumentación
de datos es la variante que obtiene el mejor resultado en validación cruzada (F1 Macro de $0.7863 \pm 0.04$), y se traslada de forma muy sólida al conjunto de test independiente, donde alcanza un F1 Macro de
0.82 y un Accuracy de 0.8. 

Tambíen se observa que la profundidad de la red es el factor mas determinante, le permite tener mas campo receptivo y resulta
en una mejor extracción de características morfológicas de la señal ECG.
Aumentar el ancho de la red o el tamaño del kernel no aporta mejoras observables.
El aumento de datos parece no aportar mejoras en este caso. Dependiendo del modelo, y de la aumentación, podría estar introduciendo
características propias de una clase en la otra.

Las puntuaciones del reto externo son menos predecibles, con un rango de $36\%-64\%$. El modelo no generaliza
a este dataset, debido a la poca representación de ese dataset en nuestro conjunto de entrenamiento completo, el modelo prioriza
la generalidad a todos los datasets, y no se especializa en el reto.

\subsection{Arquitectura SEResNet 1D con atención de canales}
\label{subsec:resultados_seresnet}

Para evaluar si los mecanismos de atención por canal (Squeeze-and-Excitation, \cite{senet1d})
ofrecen ventajas, se parametrizó la arquitectura
base y se evaluó el impacto de las variables de diseño
(profundidad, ancho, tamaño de kernel y la combinación profunda-estrecha) tanto con como sin aumento
de datos.

Los resultados obtenidos se detallan en la Tabla~\ref{tab:resultados_seresnet_variaciones}.

\begin{table}[htbp]
  \centering
  \caption{Resultados de las variaciones de la arquitectura SEResNet 1D con y sin aumento de datos}
  \label{tab:resultados_seresnet_variaciones}
  \resizebox{\textwidth}{!}{%
  \begin{tabular}{llcccc}
    \toprule
    \textbf{Configuración SEResNet 1D} & \textbf{Augment.} & \textbf{F1 Macro Val (CV)} & \textbf{F1 Macro Test} & \textbf{Accuracy Test} & \textbf{Score Reto (Event 2)} \\
    \midrule
    \multirow{2}{*}{SEResNet Base} & No & $0.7619 \pm 0.06$ & 0.80 & 0.80 & 50\% \\
                                   & Sí & $0.7426 \pm 0.05$ & 0.79 & 0.79 & 50\% \\
    \midrule
    \multirow{2}{*}{SEResNet Profunda} & No & $0.7394 \pm 0.05$ & 0.80 & 0.80 & 54\% \\
                                       & Sí & $0.7379 \pm 0.04$ & 0.78 & 0.78 & 57\% \\
    \midrule
    \multirow{2}{*}{SEResNet Ancha} & No & $0.7440 \pm 0.05$ & \textbf{0.83} & \textbf{0.83} & 54\% \\
                                    & Sí & $0.7625 \pm 0.05$ & 0.81 & 0.81 & 57\% \\
    \midrule
    \multirow{2}{*}{SEResNet Kernel Grande} & No & $0.7527 \pm 0.04$ & 0.78 & 0.78 & 50\% \\
                                            & Sí & $0.7416 \pm 0.05$ & 0.80 & 0.80 & 54\% \\
    \midrule
    \multirow{2}{*}{SEResNet Kernel Pequeño} & No & $0.7540 \pm 0.05$ & \textbf{0.83} & \textbf{0.83} & 46\% \\
                                             & Sí & $0.7583 \pm 0.02$ & 0.76 & 0.76 & 46\% \\
    \midrule
    \multirow{2}{*}{\textbf{SEResNet Profunda y Estrecha}} & No & $0.7439 \pm 0.06$ & 0.80 & 0.80 & 43\% \\
                                                           & Sí & $\mathbf{0.7680 \pm 0.05}$ & 0.79 & 0.79 & 54\% \\
    \bottomrule
  \end{tabular}%
  }
  \vskip 5pt
  \small
  \textit{Nota:} El score del reto representa la evaluación en el conjunto de test blindado de afpdb~\cite{moody_paf_2001}.
\end{table}

De los resultados de la Tabla~\ref{tab:resultados_seresnet_variaciones} concluimos que la arquitectura SEResNet Profunda y Estrecha
con aumentación de datos es la que obtiene el mejor resultado en validación cruzada (F1 Macro de $0.7680 \pm 0.05$). También obtiene
un resultado consistente en el conjunto de test independiente, con un F1 Macro de $0.79$ y un Accuracy de $0.79$.

El bloque SE permite a la red ponderar la importancia de cada canal, dependiendo de cuanta información relevante contenga cada uno.
En holters donde la señal se puede degradar por desprendimientos de electrodos o por otras razones es de especial utilidad. Se muestra
en como responde ante la aumentación de datos. Al perturbar la señal, obligamos a la red a aprender a ponderar los canales ante señales
no ideales, haciendo la red mas robusta.

Aunque obtengamos en otras variantes mejores resultados en test independiente, (las variantes SEResNet Ancha y SEResNet Kernel
Pequeño (ambas sin aumentación) obtienen un F1 de 0.83 y Accuracy del 0.83)
la SEResNet Profunda y Estrecha con aumentación de datos es la que obtiene el mejor resultado
en validación cruzada, y es la que se debe escoger de manera estricta.

El score del reto externo sigue siendo inestable, la razón es la misma que en el caso de la ResNet 1D.

\subsection{Arquitectura CNN-Transformer Híbrida y variaciones estructurales}
\label{subsec:resultados_transformer}

Para intentar superar las limitaciones del modelo base CNN-Transformer,
se comparan las siguientes variaciones estructurales entrenadas
con y sin aumento de datos:
\begin{itemize}
    \item \textbf{Transformer Base}: La arquitectura baseline con $d_{model}=128$, 4 cabezas de atención ($nhead=4$) y 3 capas de codificación.
    \item \textbf{Transformer Profundo}: Incremento del número de capas del codificador de 3 a 6 para captar relaciones temporales más abstractas.
    \item \textbf{Transformer Ancho}: Aumento de la capacidad del modelo a $d_{model}=256$, 8 cabezas ($nhead=8$) y dimensión de feedforward a 512.
    \item \textbf{Transformer Dropout Aumentado}: Elevación del factor de \textit{dropout} de 0.2 a 0.4 para mitigar el sobreajuste.
    \item \textbf{Transformer Estrecho (Narrow)}: Reducción de parámetros ($d_{model}=64$, $nhead=2$, feedforward de 128) para forzar un cuello de botella de información que prevenga la memorización.
\end{itemize}

Los resultados se detallan en la Tabla~\ref{tab:resultados_transformer_variaciones}.

\begin{table}[htbp]
  \centering
  \caption{Resultados de las variaciones de la arquitectura CNN-Transformer 1D con y sin aumento de datos}
  \label{tab:resultados_transformer_variaciones}
  \resizebox{\textwidth}{!}{%
  \begin{tabular}{llcccc}
    \toprule
    \textbf{Configuración Transformer 1D} & \textbf{Augment.} & \textbf{F1 Macro Val (CV)} & \textbf{F1 Macro Test} & \textbf{Accuracy Test} & \textbf{Score Reto (Event 2)} \\
    \midrule
    \multirow{2}{*}{Transformer Base} & No & $0.7536 \pm 0.05$ & \textbf{0.81} & \textbf{0.81} & \textbf{57\%} \\
                                      & Sí & $0.7723 \pm 0.04$ & 0.79 & 0.79 & 46\% \\
    \midrule
    \multirow{2}{*}{Transformer Profundo} & No & $0.7474 \pm 0.03$ & \textbf{0.81} & \textbf{0.81} & \textbf{57\%} \\
                                          & Sí & $0.7634 \pm 0.04$ & 0.80 & 0.80 & \textbf{57\%} \\
    \midrule
    \multirow{2}{*}{\textbf{Transformer Ancho}} & No & $0.7666 \pm 0.03$ & 0.79 & 0.79 & 46\% \\
                                       & Sí & $\mathbf{0.7727 \pm 0.05}$ & 0.80 & 0.80 & 46\% \\
    \midrule
    \multirow{2}{*}{Transformer Dropout Aumentado} & No & $0.7562 \pm 0.03$ & 0.77 & 0.78 & 46\% \\
                                                   & Sí & $0.7567 \pm 0.03$ & 0.75 & 0.75 & 46\% \\
    \midrule
    \multirow{2}{*}{Transformer Estrecho} & No & $0.7469 \pm 0.04$ & 0.78 & 0.78 & 50\% \\
                                                   & Sí & $0.7645 \pm 0.02$ & 0.77 & 0.77 & 50\% \\
    \bottomrule
  \end{tabular}%
  }
  \vskip 5pt
  \small
  \textit{Nota:} El score del reto representa la evaluación en el conjunto de test blindado de afpdb~\cite{moody_paf_2001}.
\end{table}

De los resultados de la Tabla~\ref{tab:resultados_transformer_variaciones} concluimos que la arquitectura Transformer Ancho
con aumentación de datos es la que obtiene el mejor resultado en validación cruzada (F1 Macro de $0.7727 \pm 0.05$) y con el test
independiente corroborando su generalización con un F1 Macro de 0.80 y una exactitud del 0.80.
Le sigue de cerca la variante Transformer Base con aumentación, que obtiene un F1 Macro de $0.7723 \pm 0.04$ en validación cruzada
y un F1 Macro de 0.79 en test independiente.

Disponer de mas cabezas de atención y mayor dimensión latente le permite al modelo obtener mejores resultados. En este caso
y a diferencia de las arquitecturas convolucionales, el aumento de datos si aporta mejoras en el rendimiento del modelo,
la regularización de la señal beneficia a la arquitectura de atención, que es más propensa al sobreajuste.

El score del reto externo sigue siendo inestable, la razón es la misma que en el caso de la ResNet 1D y SEResNet 1D.

\subsection{Arquitectura Mamba 1D y capacidad de representación}
\label{subsec:resultados_mamba}

El modelo Mamba 1D, basado en modelos de espacio de estados selectivos (SSM),
se evalúa como alternativa eficiente frente a las arquitecturas convolucionales y
de atención temporal. Su rendimiento se detalla en la comparativa inter-modelo
general.

\section{Evaluación comparativa inter-modelo}
\label{sec:evaluacion_cuantitativa}

El rendimiento de los mejores modelos seleccionados para cada arquitectura se evalúa a
través de las métricas: exactitud global (Accuracy), precisión, sensibilidad (Recall) y
la puntuación F1 (F1-score) tanto por clase como promedio (Macro F1).
Los resultados se detallan en la Tabla~\ref{tab:resultados_modelos}.

\begin{table}[htbp]
  \centering
  \caption{Resultados del rendimiento de las mejores variantes de cada arquitectura (seleccionadas según su F1 en validación) en validación y test}
  \label{tab:resultados_modelos}
  \resizebox{\textwidth}{!}{%
  \begin{tabular}{llcccc}
    \toprule
    \textbf{Modelo} & \textbf{Split} & \textbf{F1 Clase 0} & \textbf{F1 Clase 1} & \textbf{F1 Macro} & \textbf{Accuracy} \\
    \midrule
    \multirow{2}{*}{\textbf{ResNet 1D (Profunda, Sin Aug)}}    & Validación (CV) & --- & --- & $\mathbf{0.7863 \pm 0.04}$ & $\mathbf{0.7964 \pm 0.04}$ \\
                                                            & Test & \textbf{0.84} & \textbf{0.79} & \textbf{0.82} & \textbf{0.82} \\
    \midrule
    \multirow{2}{*}{\textbf{Transformer (Ancho, Aug)}}  & Validación (CV) & --- & --- & $0.7727 \pm 0.05$ & $0.7887 \pm 0.04$ \\
                                                            & Test & 0.82 & 0.78 & 0.80 & 0.80 \\
    \midrule
    \multirow{2}{*}{\textbf{SEResNet 1D (Profunda-Estrecha, Aug)}} & Validación (CV) & --- & --- & $0.7680 \pm 0.05$ & $0.7809 \pm 0.04$ \\
                                                            & Test & 0.82 & 0.76 & 0.79 & 0.79 \\
    \midrule
    \multirow{2}{*}{\textbf{Mamba 1D}}                      & Validación (CV) & --- & --- & $0.7138 \pm 0.09$ & $0.7350 \pm 0.09$ \\
                                                            & Test & 0.75 & 0.74 & 0.75 & 0.75 \\
    \bottomrule
  \end{tabular}%
  }
  \vskip 5pt
  \small
  \textit{Nota:} La Clase 0 corresponde a Ritmo Sinusal y la Clase 1 a la ventana crítica anterior a la Fibrilación Auricular (Pre-PAF).
\end{table}

Analizando la comparación de la Tabla~\ref{tab:resultados_modelos}, la mejor arquitectura es la ResNet 1D Profunda
sin aumentación de datos, que obtiene el mayor F1 Macro en validación cruzada ($0.7863 \pm 0.04$)
y mantiene un rendimiento sólido en el conjunto de test independiente con un F1 Macro de $0.82$ y una exactitud del $0.82$.
Tambien se observa como su F1 Clase 0 (Ritmo Sinusal) es de $0.84$ y su F1 Clase 1 (Pre-PAF) es de $0.79$,
lo que indica un buen equilibrio entre ambas clases.

Las convolucionales parecen obtener mejores resultados que las arquitecturas de atención, aunque la diferencia no es muy grande.
El modelo Mamba 1D obtiene un rendimiento sólido, con un F1 Macro de $0.7138 \pm 0.09$ en validación cruzada y un F1 Macro de $0.75$ en el conjunto de test independiente,
lo que indica una buena capacidad de generalización, aunque queda por detrás de las arquitecturas convolucionales y de atención.

\section{Análisis de resultados y discusión}
\label{sec:analisis_resultados}

%La comparación de los modelos entrenados revela conclusiones sobre el diseño de redes neuronales para señales biomédicas:

%\subsection{El impacto de la regularización estructural frente a mecanismos de atención}
%Originalmente, la superioridad del modelo \textbf{SENet 1D} sobre el baseline de \textbf{ResNet 1D} convencional se atribuía exclusivamente a la recalibración dinámica de derivaciones mediante el bloque de atención por canal (Squeeze-and-Excitation). Al permitir atenuar derivaciones ruidosas y concentrarse en las que contienen firmas biológicas claras de PAF (ondas P aberrantes, variaciones de HRV), la SEResNet base superaba al baseline original de ResNet por un margen de +7 puntos de F1 Macro.

%Sin embargo, los experimentos de optimización estructural demuestran que una regularización arquitectónica bien diseñada, como la de la red \texttt{resnet1d\_deep\_narrow\_aug} (que incrementa la profundidad a 31 capas convolucionales pero reduce drásticamente el ancho de los canales a un máximo de 192, combinada con aumentaciones de datos al vuelo), puede lograr un rendimiento equivalente (80\% F1 Macro en ambos casos). Esto indica que la constricción de capacidad estructural combinada con perturbaciones sintéticas de señal es una alternativa de regularización sumamente potente que compite directamente con la calibración adaptativa de derivaciones, logrando un rendimiento idéntico.

%\subsection{Superación del Sobreajuste en el CNN-Transformer Híbrido mediante Compresión Dimensional}
%Aunque originalmente el baseline del CNN-Transformer híbrido sufría de sobreajuste severo debido al elevado número de parámetros de los bloques de atención aplicados sobre ventanas temporales relativamente pequeñas de 10 segundos, los resultados demuestran que la compresión paramétrica resuelve esta limitación de manera drástica. Al constreñir las dimensiones latentes en la variante Transformer Estrecho ($d_{model}=64$, $nhead=2$, $dim\_feedforward=128$), se reduce la complejidad del espacio de búsqueda, lo que permite al codificador del Transformer converger con mayor facilidad y lograr una alta capacidad de generalización. Como consecuencia directa, esta variante alcanza la puntuación más alta en el reto ciego externo (71\% de Event 2 score), demostrando que la regularización por estrechez paramétrica es fundamental para aplicar arquitecturas de atención sobre señales biomédicas con conjuntos de datos limitados.

%\subsection{Modelos de Espacio de Estados (Mamba 1D) y su Capacidad de Generalización}
%La arquitectura \textbf{Mamba 1D} presenta un rendimiento sólido, obteniendo un F1 Macro del 75\% en el conjunto de test. Esto posiciona a Mamba 1D como una de las mejores arquitecturas en cuanto a generalización global, superando al baseline del CNN-Transformer Híbrido, aunque quedando por detrás de las versiones regularizadas de ResNet, SENet y del Transformer Estrecho. El éxito de Mamba 1D se puede deber a su capacidad de procesar las dependencias temporales a través de un mecanismo de selección recurrente lineal variable que escala linealmente $O(L)$ en lugar de la atención cuadrática $O(L^2)$ del Transformer. Al evitar el sobreajuste masivo que sufren los Transformers en conjuntos de datos medianos, Mamba conserva gran parte de la información de largo alcance sin penalizar su capacidad de generalización en test.

\subsection{Discrepancia entre el rendimiento en test general y el benchmark de PhysioNet}
Un hallazgo clave de este trabajo es la discrepancia observada entre el rendimiento en el
conjunto de prueba general y las puntuaciones en el reto oficial de PhysioNet (\texttt{afpdb}):
\begin{itemize}
  \item Originalmente se fue más estricto a la hora de extraer ventanas, asumiendo que si no se
  encontraban 5 minutos completos, no se consideraba. Esto dejó a la base de datos \texttt{afpdb} como una fuente
  altamente representativa, y por ello los modelos obtenían puntuaciones altas en el reto (hasta 79\% en el
  CNN-Transformer y 75\% en la ResNet 1D), pero a costa de un F1-score muy bajo en condiciones generales
  (56\% y 61\% respectivamente).
  \item Tras realizar el aumento de datos mediante ventanas variables, el volumen de la base de datos
  del reto (\texttt{afpdb}) pasa a representar solo el 2.65\% del total de Clase 1 en el entrenamiento,
  que ahora es dominado por registros de Holter ambulatorio de 24 horas (\texttt{ltafdb}) y
  dispositivos corporales inteligentes (\texttt{shdb-af}).
  \item Como consecuencia directa, la generalización ante nuevos sujetos
  mejora notablemente (alcanzando el 82\% de F1 en la ResNet optimizada),
  pero las puntuaciones del reto específico se vuelven inestables.
\end{itemize}

Al optimizar los modelos con un dataset más grande, las redes aprenden a priorizar rasgos
generales de ECG ambulatorio y wearables. Sin embargo, esto diluye la especialización del
modelo a la base de datos \texttt{afpdb} original.

\subsection{Estudio de la contribución de la variabilidad del ritmo cardíaco (HRV)}
\label{subsec:estudio_ablacion_hrv}

Para evaluar el impacto de las variables de variabilidad del ritmo cardíaco (HRV) en el
proceso de decisión del clasificador, se comparó un modelo (el híbrido \textbf{Transformer Estrecho con
Aumento de Datos}) frente a dos configuraciones:
\begin{itemize}
    \item \textbf{Configuración de solo señal}: Se reentrenó la misma arquitectura del
    Transformer Estrecho y Aumentado eliminando por completo la rama de entrada de HRV,
    forzando a la red a basar sus predicciones únicamente en los 60 segundos de morfología de
    señal de ECG.
    \item \textbf{Configuración de solo características HRV}: Se entrenaron dos clasificadores clásicos
    (un \textit{Random Forest}, y un \textit{Gradient Boosting})
    que recibieron exclusivamente las
    9 características de HRV locales calculadas a partir de las ventanas de ritmo sinusal,
    aisladas de cualquier información morfológica del ECG.
\end{itemize}

Los resultados obtenidos mediante esta estrategia se presentan en la Tabla~\ref{tab:resultados_ablacion_hrv}.

\begin{table}[htbp]
  \centering
  \caption{Resultados del estudio de contribución de HRV frente al modelo híbrido completo}
  \label{tab:resultados_ablacion_hrv}
  \resizebox{\textwidth}{!}{%
  \begin{tabular}{llcccc}
    \toprule
    \textbf{Modelo} & \textbf{Modalidad de Entrada} & \textbf{F1 Macro Val (CV)} & \textbf{F1 Macro Test} & \textbf{Accuracy Test} & \textbf{Score Reto (Event 2)} \\
    \midrule
    \textbf{Transformer Narrow (Aug)} & Híbrido (ECG + HRV) & $\mathbf{0.7645 \pm 0.02}$ & 0.77 & 0.77 & 50\% \\
    Transformer Narrow (Aug) & Solo ECG & $0.7518 \pm 0.02$ & 0.76 & 0.76 & 50\% \\
    \midrule
    Random Forest  & Solo HRV  & $0.7137 \pm 0.05$ & \textbf{0.82} & \textbf{0.82} & \textbf{71\%} \\
    Gradient Boosting  & Solo HRV  & $0.6978 \pm 0.04$ & 0.82 & 0.82 & 61\% \\
    \bottomrule
  \end{tabular}%
  }
  \vskip 5pt
  \small
  \textit{Nota:} El score del reto representa la evaluación en el conjunto de test blindado de afpdb~\cite{moody_paf_2001}.
\end{table}

De los resultados de la Tabla~\ref{tab:resultados_ablacion_hrv} se extraen varias conclusiones.
El objetivo del experimento era evaluar si el análisis de las redes neuronales a partir
de $60\text{ s}$ de ECG es competitivo frente a modelos clásicos basados únicamente en variables de
HRV, y si la combinación de ambas fuentes de
información mejora el clasificador.

En primer lugar, el modelo híbrido (\textit{ECG + HRV}) supera al
modelo basado en \textit{Solo ECG} en validación ($0.7645$ frente a $0.7518$) y en test local ($0.77$ frente a $0.76$ de F1 Macro).
Lo que nos dice que la red se beneficia de la información que le aportan las variables HRV.

El clasificador basado únicamente en características HRV (\textit{Random Forest}) obtiene un F1 Macro de $0.7137 \pm 0.05$
en validación cruzada, y un F1 Macro de $0.82$ en el conjunto de test independiente, lo que indica que las variables HRV contienen
información de gran utilidad para la predicción de fibrilación auricular. Aunque no obtenga en validación el mejor resultado, no
se queda atrás.

En concreto obtiene un $71\%$ de acierto en el reto PhysioNet. Dado que el volumen de datos clínicos
del estudio es acotado, las redes neuronales, por muy complejas que sean, no disponen de datos suficientes
para aprovechar su capacidad de analizar la morfología de la señal, lo que justifica la competitividad del
modelo clásico.


En segundo lugar, al evaluar la generalización clínica en el test externo del reto, aunque el modelo clasico
se posicionaba como el mejor, no obtiene resultados competitivos, debido tanto a la inestabilidad que hemos
comentado a lo largo del trabajo del que presenta el score del reto, como quiza la falta de generalización
de un modelo que no tiene en cuenta la morfología de la señal que se está tratando.


\section{Comparación con trabajos relacionados}
\label{sec:comparacion_trabajos}

La mayoría de los algoritmos publicados en la literatura científica que evalúan su rendimiento en el
PAF Prediction Challenge (afpdb\cite{moody_paf_2001}) entrenan sus modelos únicamente utilizando
los datos de la base de datos \texttt{afpdb}, o de bases de datos específicas.
Aunque algunos logran clasificaciones sobresalientes en este benchmark específico (79\% de precisión),
carecen de la generalidad que ofrecen estos modelos.

Este trabajo demuestra que el entrenamiento exclusivo sobre una única base de datos induce un sesgo de
adquisición artificial. Al contrastar nuestros modelos con una validación cruzada inter-sujeto
sobre cinco bases de datos, demostramos la viabilidad de un predictor de PAF generalizable capaz de
mantener una alta sensibilidad y especificidad incluso en señales capturadas por wearables en
condiciones cotidianas.
